\section{\ref{sec:serocomputer}장. \nt{SERO} 뉴런}

\nt{SERO} 뉴런 자체의 성질(리듬, 자가 억제, 수용체에 따른 부호, 하위 시스템)은 직접 측정되었다. 이 장의 계산, 즉 조절기, 필터, 논리는 이 장의 식에서 나온다. $\tau_c$, 세로토닌의 확산 계수, 방출 지점의 수는 추정이다. 고리가 뇌에서 수준 조절기로 작동한다는 것은 가설이고(봉선핵에서 자가 억제는 점대점이며 짧은 멈춤만 만든다), \nt{SERO}가 지평을 정한다는 것은 이를 지지하는 행동 실험이 있는 가설이다.

{\footnotesize
\setlength{\tabcolsep}{3pt}\renewcommand{\arraystretch}{1.15}
\begin{longtable}{>{\raggedright}p{0.5cm}>{\raggedright}p{4.0cm}>{\raggedright}p{5.6cm}>{\raggedright}p{2.3cm}>{\raggedright\arraybackslash}p{1.7cm}}
\toprule
No. & 명제 & 실험과 측정 내용 & 출처 & 상태 \\
\midrule\endfirsthead
\toprule
No. & 명제 & 실험과 측정 내용 & 출처 & 상태 \\
\midrule\endhead
1 & 세로토닌 작용의 부호는 수신자의 수용체가 고른다(명제~\ref{prop:receptor}) & 세로토닌 수용체는 약리와 기전에 따라 여러 계열로 나뉜다. 5-HT$_{1A}$는 G 단백질을 통해 칼륨 채널을 열어 세포를 억제하고, 5-HT$_{2A}$와 이온성 5-HT$_3$는 흥분시킨다. 한 신경전달물질이 세포에 따라 반대 응답을 낸다. & \cite{BarnesSharp1999} & 측정됨 \\
2 & 깨어 있을 때 \nt{SERO} 뉴런은 초당 몇 개의 스파이크를 고르게 낸다 & 자유롭게 움직이는 고양이의 등쪽 봉선핵 뉴런 기록: 조용히 깨어 있을 때 $2.8\pm0.2$ Hz의 느리고 규칙적인 발화. 서파 수면에서는 발화율이 $34$--$68\,\%$, 렘수면에서는 $98\,\%$ 떨어진다. 마취한 쥐에서는 $1.7\pm0.5$ Hz로 매우 규칙적이다. & \cite{TrulsonJacobs1979, GagnonParent2014, JacobsAzmitia1992} & 측정됨 \\
3 & \nt{SERO} 뉴런은 페이스메이커다. 스파이크마다 밀어 줄 필요는 없지만 고른 배경 입력은 필요하다(절편에서는 $\alpha_1$ 수용체를 통한 노르아드레날린, 생체에서는 \nt{NORA}에서 온다) & 뇌 절편에서 세포는 느리고 고르게 발화하며, 스파이크마다 $200$--$800\ms$의 후과분극이 뒤따른다. 절편에서 자발적으로 발화하는 세포는 생체보다 적다. 고른 리듬에는 $\alpha_1$ 수용체를 통한 노르아드레날린의 지속적 입력이 필요하고, 이를 차단하면 리듬이 꺼진다. & \cite{VandermaelenAghajanian1983} & 측정됨 \\
4 & 수용체는 거의 모두 대사성이다. 응답은 느려서 약 1초 안에 오르고 내린다 & 5-HT$_3$을 뺀 모든 수용체 계열은 G 단백질과 결합한다. 절편에서 유발된 세로토닌 방출은 5-HT$_{1A}$를 통해 1초 안에 오르고 내리는 억제 전류를 만든다. & \cite{BarnesSharp1999, CourtneyFord2016} & 측정됨 \\
5 & 목표 영역에서 세로토닌은 틈뿐 아니라 부피로도 나온다. 봉선핵 자체에서 이웃 억제는 점대점이다 & 절편에서의 고속 주사 순환 전압전류법: 스파이크 하나당 방출량이 시냅스가 있는 영역과 시냅스가 거의 없는 봉선핵에서 같다. 재흡수와 수용체를 차단해도 달라지지 않는다. 말단당 분자 수가 수송체와 수용체보다 훨씬 적고, 세포 밖 농도는 주 수용체의 친화도와 일치한다. 봉선핵에서 5-HT$_{1A}$를 통한 억제는 점대점이다. 수송체가 세로토닌이 틈 밖에 쌓이지 못하게 한다. & \cite{Bunin1998, CourtneyFord2016} & 측정됨 \\
6 & 식~\eqref{eq:seroneuron}의 농도는 재흡수로 줄어든다. $\tau_c$가 1초 정도라는 것은 추정이다 & 생체 뇌에서의 전압전류법: 세포 밖 세로토닌을 제한하는 것은 합성이 아니라 주로 재흡수와 분해다. 인용한 연구에 $\tau_c$의 직접 측정은 없다. 크기 정도는 이 장의 추정이다. & \cite{Hashemi2012serotonin} & 측정됨; $\tau_c$는 추정 \\
7 & 페이스메이커 하나로 만드는 NOT과 NOR; \nt{SERO} 논리의 완전성(명제~\ref{prop:serocomplete}, 그림~\ref{fig:serogates}) & 식~\eqref{eq:seroneuron}에서 나온다. 억제되는 페이스메이커는 NOR이고, NOR은 함수적으로 완전하다. \nt{SERO} 뉴런으로 논리를 만든 실험은 없다. 부피가 분리되어 있어야 한다는 조건은 뇌에서 성립하지 않는다. & 이 장의 유도 & 이론 \\
8 & 세로토닌은 \nt{SERO} 뉴런 자신에 있는 수용체를 통해 그 뉴런을 억제한다 & 마취한 쥐에서 5-HT$_{1A}$ 작용제를 정맥 주사하거나 이온영동으로 주면, 세로토닌 자체와 같은 용량--반응 관계로 등쪽 봉선 뉴런의 발화가 억제된다. 5-HT$_{1B}$ 작용제는 거의 효과가 없다. 절편에서 이웃 세포의 내인성 세로토닌은 5-HT$_{1A}$를 통해 전류와 발화의 짧은 멈춤을 만들지만 평균 발화율은 바꾸지 않는다. & \cite{Sprouse1987, CourtneyFord2016} & 측정됨 \\
9 & 모델에서 공통 부피를 통한 고리는 수준 조절기다. 편차와 시간 상수가 $1+G$배 줄어든다~\eqref{eq:feedback}. 고리가 뇌에서 이렇게 작동한다는 것은 가설이다 & 음의 되먹임 증폭기의 고전적 공식이며, 이 장에서 새로 유도했다. \nt{SERO} 시스템의 고리 이득 $G$는 측정되지 않았다. 측정된 것은 절편에서 자가 억제가 점대점이고, 짧은 멈춤을 만들며, 평균 발화율을 바꾸지 않는다는 점이다. & \cite{Black1934feedback, CourtneyFord2016} & 이론; 뇌에서는 가설 \\
10 & 농도는 발화율의 이동 평균, 즉 저역 통과 필터다~\eqref{eq:lowpass}, 그림~\ref{fig:lowpass} & 선형 방정식~\eqref{eq:seroneuron}의 해. 그림은 $\tau_c=1\,\text{s}$에서 이 공식으로 계산한 것이다. \texttt{models/}에 별도 모델은 없다. & 이 장의 유도 & 이론 \\
11 & 틈의 굴곡은 작은 분자의 확산을 약 $2.6$배 늦춘다 & 점원 방법: 테트라메틸암모늄을 이온영동으로 주입하고 이온 선택성 전극으로 그 농도를 절편과 생체 뇌에서 기록한다. 굴곡도 $\lambda\approx1.6$, 즉 유효 $D$는 자유 확산보다 $\lambda^2\approx2.6$배 작다. 세포 사이 부피의 비율은 약 $0.2$. 이 연구들에서 조직 속 세로토닌 자체의 확산 계수는 측정하지 않았다. 이 장에서 그 값은 추정이다. & \cite{Sykova2008} & 측정됨 \\
12 & 독립 “전선”의 길이 $\ell\sim\sqrt{D\tau}\approx20$~\textmu m~\eqref{eq:diffusionlength}; 전선의 수는 이런 영역의 수로 제한된다(명제~\ref{prop:volumewires}) & 확산 법칙에 $D$와 $\tau$의 추정값(6행과 11행)을 대입한 것이며, 명제~\ref{prop:volumewires}는 그 따름정리다. 체적 전달의 독립 채널 수를 센 직접 실험은 없다. & 이 장의 유도 & 이론 \\
13 & \nt{SERO} 뉴런 하나의 출력은 뇌의 넓은 영역에 퍼져 있다. 방출 지점은 추정으로 $\sim10^5$개 & 쥐 등쪽 봉선 뉴런 하나하나의 완전 재구성: 모든 상행 축삭이 많이 갈라지며, 세포 하나의 축삭 총길이는 최대 $18.7$~cm, 한 세포의 가지가 선조체, 전전두피질, 편도체에 닿는다. 세포당 방출 지점 수는 직접 세지 않았다. & \cite{GagnonParent2014} & 측정됨; $10^5$은 추정 \\
14 & \nt{SERO} 뉴런은 출력과 응답이 다른 몇 개의 하위 시스템으로 나뉜다 & 생쥐에서의 바이러스 유전 표지: 편도체로 가는 뉴런과 전두피질로 가는 뉴런은 분지가 거의 겹치지 않고, 서로 다른 입력을 받으며, 불쾌한 자극에 반대로 반응한다. 각 집단을 켜고 끄면 행동이 서로 다르게 바뀐다. & \cite{Ren2018} & 측정됨 \\
15 & 지수형 할인에서 인내 조건은 $D<\tau_\gamma\ln(R/r_0)$~\eqref{eq:patience}; 측정된 쌍곡선형 할인에서는 $D<\tau_\gamma(R/r_0-1)$ & 할인 $e^{-D/\tau_\gamma}$ 또는 $1/(1+D/\tau_\gamma)$에서 나오며, 이 장에서 유도했다. 몇 초 지연에서 선택하는 원숭이: 할인은 지수보다 쌍곡선에 가깝다. 지연 보상을 알리는 신호에 대한 \nt{DOPA} 뉴런의 응답도 지연에 따라 쌍곡선으로 떨어진다. & 이 장의 유도; \cite{KobayashiSchultz2008} & 이론 \\
16 & \nt{SERO}가 지평 $\tau_\gamma$를 정한다(\ref{sec:horizon}절) & 생쥐에서 등쪽 봉선 \nt{SERO} 뉴런을 광유전학으로 켜면 지연 보상을 더 오래 기다리고, 드물어진 보상을 더 끈질기게 찾는다. 사람에서 세로토닌을 낮추면(트립토판을 뺀 식단) 지연 보상의 할인이 가팔라진다. 농도가 어떻게 $\tau_\gamma$로 바뀌는지는 알려져 있지 않다. & \cite{Doya2002, Miyazaki2014, Lottem2018, Schweighofer2008} & 가설 \\
\bottomrule
\end{longtable}}
